\section{Replicación de sistemas complejos naturales}\label{sec:exp2}

En este experimento las evoluciones de entrenamiento no provienen de una regla conocida, sino de observaciones de sistemas reales: un incendio forestal visto desde un dron (sección \ref{sec:exp2_flame3}), el desarrollo de embriones en una incubadora con cámara (sección \ref{sec:exp2_embriones}) y el movimiento y división de células en cultivo (sección \ref{sec:exp2_ctc}). Los tres comparten las características que permiten modelarlos como un autómata celular: se observan como una malla bidimensional con encuadre fijo, el estado de cada región depende principalmente de su vecindad y existen secuencias con un paso de tiempo regular.

A diferencia del experimento anterior, los datos no están en el formato del simulador. Cada sistema requiere un método de limpieza propio que lleve las imágenes a una evolución $(T, n, n)$ de estados discretos con significado físico, con dos criterios comunes: los umbrales que definen los estados se calculan con toda la secuencia y no por imagen, para que un estado signifique lo mismo en todo momento; y la malla es la misma para todos los pasos de tiempo, para que el autómata no confunda un cambio de encuadre con una dinámica.
